%

\chapter{Aksjomatyka}

Tekst rozdziału Aksjomatyka. \loremipsum
Omówimy w tej sekcji najstarszy układ aksjomatów dla geometrii (płaskiej).
Wskażemy, jakie błędy popełni Euklides, ale też jak zostaną naprawione przez Hilberta i jego kolegów.
Kilka słów o polu powierzchni.
Pierwsze i ostatnie (?) wystąpienie origami.

\todofoot{https://www.deltami.edu.pl/2015/01/prozny-trud/}

\input{chapters/01-aksjomatyka/euklides}
\input{chapters/01-aksjomatyka/hilbert}
%

\section{Aksjomaty Tarskiego}
\index{aksjomat!Tarskiego}%

Zmienne są punktami, relacja $\beta$ mówi, że punkty leżą na prostej w kolejności takiej, jak w napisie, zaś relacja $\equiv$ mówi, że dwa odcinki są przystające:
\begin{enumerate}
    \item jeśli $\beta (x, y, x)$, to $x = y$;
    \item jeśli $\beta (x, y, u)$ oraz $\beta (y, z, u)$,  to $\beta (x, z, u)$;
    \item jeśli $\beta (x, y, z)$, $\beta (x, y, u)$ oraz $x \neq y$, to $\beta (x, z, u)$ lub $\beta (x, u, z)$;
    \item $xy \equiv yx$;
    \item jeśli $xy \equiv zu$ oraz $xy \equiv vw$, to $zu \equiv vw$;
    \item jeśli $xy \equiv zz$, to $x = y$;
    \item aksjomat Pascha: jeśli $\beta (x, t, u)$, $\beta (y, u, z)$, to istnieje taki $v$, że $\beta (x, v, y)$, $\beta(z, t, v)$'';
    \item V aksjomat Euklidesa\footnote{Ten aksjomat można zastąpić przez ,,na każdym trójkącie można opisać okrąg'' (mniej zmiennych) albo ,,linia środkowa w trójkącie jest dwukrotnie krótsza od podstawy'' (brak kwantyfikatorów).} (dla każdego kąta i punktu wewnątrz istnieje odcinek przez ten punkt, który przecina ramiona kąta): jeśli $\beta(x, u, t)$, $\beta (y, u, z)$, $x \neq y$, to istnieją takie $v, w$, że $\beta(x, z, v)$, $\beta(x, y, w)$ oraz $\beta (v, t, w)$;
    \item aksjomat pięciu odcinków (cecha przystawania bkb): jeśli $xy \equiv x'y'$, $yz \equiv y'z'$, $xu \equiv x'u'$, $yu \equiv y'u'$, $\beta (x, y, z)$, $\beta(x', y', z')$ i $x \neq y$, to $zu \equiv z'u'$;
    \item odkładanie odcinków: istnieje $z$ taki, że $\beta (x, y, z)$ i $yz \equiv uw$;
    \item istnieją trzy niewspółliniowe punkty\footnote{Bez tego aksjomatu teoria ma model na prostej rzeczywistej.} $x, y, z$: $\neg \beta(x, y, z)$, $\neg \beta(y, z, x)$, $\neg \beta(z, x, y, )$;
    \item trzy punkty równoodległe od końców odcinka tworzą prostą\footnote{Bez tego aksjomatu teoria ma model w trzech i więcej wymiarach.}: jeśli $xu \equiv xv$, $yu \equiv yv$, $zu \equiv zv$ i $u \neq v$, to $\beta(x, y, z)$, $\beta(y, z, x)$ lub $\beta(z, x, y)$;
    \item schemat aksjomatu ciągłości.
\end{enumerate}

% TODO: https://en.wikipedia.org/wiki/Tarski%27s_axioms

% TODO: https://smp.uws.edu.pl/msn/14/10-18.pdf
% Kompletne, beznadziejne bezdroża, Marek KORDOS

% https://www.deltami.edu.pl/media/articles/2015/01/delta-2015-01-prozny-trud.pdf

Patrz też Greenberg \cite{greenberg_2010}. % TODO: czemu patrz?

%
%

\section{Aksjomaty Birkhoffa} % lang-pl
\section{Assiomi di Birkhoff} % lang-it
\index{aksjomat!Birkhoffa}% % lang-pl
\index{assioma!di Birkhoff}% % lang-it
Aksjomaty Birkhoffa z 1932 roku. % lang-pl
Assiomi di Birkhoff del 1932. % lang-it

% TODO: https://en.wikipedia.org/wiki/Birkhoff%27s_axioms

%
\input{chapters/01-aksjomatyka/neutral}
% eksperymentalne tłumaczenie na włoski

%

\section{Elementarne wyniki} % lang-pl
\section{Risultati elementari} % lang-it
\begin{definition}[asse del segmento] % lang-it
    \label{def_symetralna} % lang-it
    La retta perpendicolare a un segmento e passante per il suo punto medio si chiama asse del segmento. % lang-it
    \index{asse del segmento} % lang-it
\end{definition} % lang-it


\begin{definition}[symetralna] % lang-pl
    \label{def_symetralna} % lang-pl
	Prostą prostopadłą do odcinka i przechodzącą przez jego środek nazywamy symetralną. % lang-pl
    \index{symetralna} % lang-pl
\end{definition} % lang-pl
Symetralna odcinka $AB$ jest miejscem geometrycznym punktów, które są równodległe od obu końców odcinka, a zatem są środkiem pewnego okręgu przechodzącego przez punkty $A$ oraz $B$ (Guzicki \cite[s. 14, 15]{guzicki_2021}). % lang-pl
L'asse del segmento $AB$ è il luogo geometrico dei punti equidistanti dai due estremi del segmento e quindi costituisce l'insieme dei centri delle circonferenze passanti per i punti $A$ e $B$ (Guzicki \cite[s. 14, 15]{guzicki_2021}). % lang-it


W szczególności, trzy współliniowe punkty nie leżą na jednym okręgu, ale już trzy niewspółliniowe punkty wyznaczają taki okrąg jednoznacznie. % lang-pl
In particolare, tre punti allineati non appartengono a una stessa circonferenza, mentre tre punti non allineati determinano univocamente una tale circonferenza. % lang-it


\index{współliniowy} % lang-pl
\index{allineato} % lang-it


Okrąg przechodzący przez wszystkie wierzchołki wielokąta nazywamy opisanym na nim, zatem poprzednie stwierdzenie można wysłowić krótko ,,na każdym trójkącie można opisać dokładnie jeden okrąg''. % lang-pl
La circonferenza passante per tutti i vertici di un poligono si dice circoscritta al poligono; pertanto l'affermazione precedente può essere formulata brevemente come ``per ogni triangolo esiste una e una sola circonferenza circoscritta''. % lang-it


Mówimy też, że wielokąt jest wpisany w okrąg. % lang-pl
Si dice anche che il poligono è inscritto nella circonferenza. % lang-it


\index{okrąg!opisany} % lang-pl
\index{circonferenza!circoscritta} % lang-it

% TODO: https://en.wikipedia.org/wiki/Bisection#Line_segment_bisector
% TODO: https://en.wikipedia.org/wiki/Circumcircle

\begin{proposition}
	\label{hartshorne_52x}
    Niech $AB$ będzie odcinkiem. % lang-pl
    
    Sia $AB$ un segmento. % lang-it

	Istnieje wtedy trójkąt równoramienny, którego podstawą jest $AB$. % lang-pl
    Allora esiste un triangolo isoscele avente $AB$ come base. % lang-it
    
    \index{trójkąt!równoramienny} % lang-pl
    \index{triangolo!isoscele} % lang-it
\end{proposition}

Powyższe stwierdzenie jest ciekawe, bo jest prawdziwe na płaszczyźnie Hilberta: jego prawdziwość nie zależy od aksjomatu Pascha. % lang-pl
L'affermazione precedente è interessante perché è valida nel piano di Hilbert: la sua validità non dipende dall'assioma di Pasch. % lang-it


\index{płaszczyzna!Hilberta} % lang-pl
\index{piano!di Hilbert} % lang-it


\index{aksjomat!Pascha} % lang-pl
\index{assioma!di Pasch} % lang-it


W geometrii nieeuklidesowej może nie istnieć trójkąt równoboczny o danej podstawie. % lang-pl
Nella geometria non euclidea può non esistere un triangolo equilatero con una base assegnata. % lang-it


Ale na płaszczyźnie Hilberta z aksjomatem (P), bez zakładania aksjomatu (E), istnieje trójkąt równoboczny o danej podstawie -- patrz ćwiczenie u Hartshorne'a \cite[s. 184]{hartshorne2000}. % lang-pl
Tuttavia, nel piano di Hilbert con l'assioma (P), senza assumere l'assioma (E), esiste un triangolo equilatero con una base assegnata; si veda l'esercizio in Hartshorne \cite[s. 184]{hartshorne2000}. % lang-it

\begin{proposition}
\label{hartshorne_52}%
	Linia środkowa (odcinek łączący środki pewnych dwóch boków trójkąta) jest równoległa do trzeciego boku. % lang-pl
    Il segmento che congiunge i punti medi di due lati di un triangolo è parallelo al terzo lato. % lang-it
    
    \index{segmento dei punti medi} % lang-it

    \index{linia środkowa} % lang-pl
    
    Jej długość jest dwukrotnie mniejsza od długości tego boku. % lang-pl
    La sua lunghezza è pari alla metà della lunghezza di tale lato. % lang-it
\end{proposition}
% Hartshorne s. 52

Tego stwierdzenia nie ma w Elementach Euklidesa, ale można wyprowadzić je z księgi I (I.29, I.26, I.34), jak wspomina Hartshorne \cite[s. 45, 52. 53]{hartshorne2000}. % lang-pl
Questa affermazione non compare negli \emph{Elementi} di Euclide, ma può essere dedotta dal Libro I (I.29, I.26, I.34), come osserva Hartshorne \cite[s. 45, 52--53]{hartshorne2000}. % lang-it


Patrz też Bogdańska, Neugebauer \cite[s. 15, 52]{neugebauer_2018}. % lang-pl
Si vedano anche Bogdańska e Neugebauer \cite[s. 15, 52]{neugebauer_2018}. % lang-it


Joanna Jaszuńska napisze ,,Niby nic'' w $\Delta_{17}^5$. % lang-pl
Joanna Jaszuńska pubblicherà l'articolo ``Niby nic'' in $\Delta_{17}^5$. % lang-it

% \todofoot{https://www.deltami.edu.pl/2017/05/niby-nic/}

\begin{corollary}
	Niech $ABC$ będzie trójkątem, zaś punkty $D$, $E$ i $F$ środkami jego boków. % lang-pl
    Sia $ABC$ un triangolo e siano $D$, $E$ e $F$ i punti medi dei suoi lati. % lang-it

	Wtedy cztery małe trójkąty utworzone na bokach $DE$, $EF$, $FD$ są przystające do siebie. % lang-pl
    Allora i quattro piccoli triangoli determinati dai segmenti $DE$, $EF$ e $FD$ sono congruenti tra loro. % lang-it
\end{corollary}

Do tego wniosku potrzeba dodatkowo cechy przystawania bok-bok-bok (I.8). % lang-pl
Per ottenere questa conclusione è inoltre necessario utilizzare il criterio di congruenza lato-lato-lato (I.8). % lang-it


\index{cecha przystawania!bok-bok-bok} % lang-pl
\index{criterio di congruenza!lato-lato-lato} % lang-it


Pompe \cite[s. 40]{pompe_2022} uogólnia fakt \ref{hartshorne_52} do czworokątów: % lang-pl
Pompe \cite[s. 40]{pompe_2022} generalizza il fatto \ref{hartshorne_52} ai quadrilateri: % lang-it

\begin{proposition}
	Dany jest czworokąt wypukły $ABCD$ oraz punkty $K$ i $L$ leżące na środkach boków $AD$ i $BC$. % lang-pl
    Siano dati un quadrilatero convesso $ABCD$ e i punti $K$ e $L$, punti medi rispettivamente dei lati $AD$ e $BC$. % lang-it

	Wtedy $KL \le \frac 12 (AB +CD)$ z równością wtedy i tylko wtedy, gdy proste $AB$ i $CD$ są równoległe. % lang-pl
    Allora $KL \le \frac 12 (AB + CD)$, con uguaglianza se e solo se le rette $AB$ e $CD$ sono parallele. % lang-it

	(Wtedy prosta $KL$ też jest równoległa do nich). % lang-pl
    (In tal caso anche la retta $KL$ è parallela ad esse). % lang-it
\end{proposition}

\begin{corollary}
	Prosta przechodząca przez środek ramienia trapezu i równoległa do jego podstaw przechodzi przez środek drugiego ramienia oraz środki obu przekątnych. % lang-pl
    La retta passante per il punto medio di un lato obliquo di un trapezio e parallela alle sue basi passa per il punto medio dell'altro lato obliquo e per i punti medi di entrambe le diagonali. % lang-it
\end{corollary}

\begin{definition}[dwusieczna] % lang-pl
	Niech $\angle ABC$ będzie kątem. % lang-pl
\begin{definition}[bisettrice] % lang-it
    Sia $\angle ABC$ un angolo. % lang-it

	Półprostą $BD$ leżącą w jego wnętrzu taką, że kąty $\angle ABD = \angle DBC$ mają równe miary, nazywamy dwusieczną kąta $\angle ABC$. % lang-pl
    La semiretta $BD$, contenuta nell'interno dell'angolo e tale che gli angoli $\angle ABD$ e $\angle DBC$ abbiano la stessa ampiezza, si chiama bisettrice dell'angolo $\angle ABC$. % lang-it
\end{definition}

Dwusieczna kąta (wypukłego) jest miejscem geometrycznym punktów, które leżą w jego wnętrzu i~są równoodległe od ramion. % lang-pl
La bisettrice di un angolo (convesso) è il luogo geometrico dei punti che giacciono nel suo interno e sono equidistanti dai suoi lati. % lang-it


Zadania o niej znajdą się w $\Delta_{16}^{4}$. % lang-pl
Esercizi su questo argomento compariranno in $\Delta_{16}^{4}$. % lang-it


Symetralne i dwusieczne jako bliźnięta jedno- bądź dwujajowe będą tematem artykułu w $\Delta_{13}^8$. % lang-pl
Gli assi dei segmenti e le bisettrici, considerate come gemelli monozigoti o dizigoti, saranno il tema di un articolo in $\Delta_{13}^8$. % lang-it

% \todofoot{https://www.deltami.edu.pl/2016/04/dwusieczne/}
% \todofoot{https://www.deltami.edu.pl/2013/08/symetralna-i-dwusieczna-bliznieta-jedno-czy-dwujajowe} %  + www.deltami.edu.pl/media/issues/2013/08/delta-2013-08.pdf

\begin{proposition}
	Dwusieczna $t$ kąta przy podstawie trójkąta równoramiennego o podstawie $a$ oraz ramionach $b, b$ spełnia podwójną nierówność % lang-pl
	
	La bisettrice $t$ di un angolo alla base di un triangolo isoscele di base $a$ e lati congruenti $b, b$ soddisfa la doppia disuguaglianza % lang-it
	\begin{equation}
		\sqrt{2} \cdot \frac{ab}{a+b} < t < 2 \cdot \frac{ab}{a+b}.
	\end{equation}
\end{proposition}
% https://en.wikipedia.org/wiki/List_of_triangle_inequalities#Isosceles_triangle

%
% eksperymentalne tłumaczenie na włoski

\section{Pole powierzchni} % lang-pl
\section{Area} % lang-it
\label{sekcja_pole_powierzchni}

Pole powierzchni opisuje, jak duża jest dana figura płaska. % lang-pl
L'area descrive quanto è grande una data figura piana. % lang-it
Opiszemy, jak na przestrzeni wieków zmieniać się będzie matematyczny opis pola, skupiając się na rozcięciach. % lang-pl
Descriveremo come, nel corso dei secoli, cambierà la descrizione matematica dell'area, concentrandoci sulle dissezioni. % lang-it
Sekcję zwieńczy twierdzenie Wallace'a-Bolyaia-Gerwiena, charakteryzujące figury o równych polach. % lang-pl
La sezione si concluderà con il teorema di Wallace-Bolyai-Gerwien, che caratterizza le figure di area uguale. % lang-it
Miarę Lebesgue'a wzmiankujemy krótko na koniec w formie ciekawostki, nie będzie nam do niczego potrzebna. % lang-pl
Accenneremo brevemente alla misura di Lebesgue alla fine, come curiosità: non ci servirà a nulla. % lang-it

Euklides nawet nie zdefiniuje on figur o równych polach, a tym bardziej samego pola jako liczby. % lang-pl
Euclide non definirà nemmeno le figure di area uguale, tanto meno l'area stessa come numero. % lang-it
(Będzie unikać problemu kwadratury koła, ale w księdze XII napisze, że pola dwóch kół mają się do siebie tak, jak kwadraty ich promieni. % lang-pl
(Eviterà il problema della quadratura del cerchio, ma nel libro XII scriverà che le aree di due cerchi stanno tra loro come i quadrati dei loro raggi. % lang-it
Czytaj dalej w sekcji \ref{sekcja_przyblizenia_pi}.) % lang-pl
Continua a leggere nella sezione \ref{sekcja_przyblizenia_pi}.) % lang-it
Ale czytając jego dowody możemy domyślić się, że równopolowe figury mają następujące właściwości: % lang-pl
Ma leggendo le sue dimostrazioni possiamo intuire che le figure equivalenti (equal-area) hanno le seguenti proprietà: % lang-it
\begin{enumerate}
\item przystające figury mają równe pola; % lang-pl
\item le figure congruenti hanno aree uguali; % lang-it
\item suma, różnica i połówki figur o równych polach mają równe pola; % lang-pl
\item la somma, la differenza e le metà di figure di area uguale hanno aree uguali; % lang-it
\item całość jest większa od części i wreszcie % lang-pl
\item il tutto è maggiore della parte e, infine, % lang-it
\item boki kwadratów o równych polach są równej długości. % lang-pl
\item i lati di quadrati di area uguale hanno la stessa lunghezza. % lang-it
\end{enumerate}

Czwarty punkt jest wnioskiem z trzeciego. % lang-pl
Il quarto punto è una conseguenza del terzo. % lang-it
Hartshorne \cite[s. 196--204]{hartshorne_2000} podczas formalizacji pojęcia pola dla płaszczyzny Hilberta wspomni aksjomat de Zolta: % lang-pl
Hartshorne \cite[p. 196--204]{hartshorne_2000}, formalizzando il concetto di area per il piano di Hilbert, menzionerà l'assioma di De Zolt: % lang-it

\begin{axiom}
[Z, de Zolta] % lang-pl
[Z, di De Zolt] % lang-it
\index{aksjomat!de Zolta}% % lang-pl
\index{assioma!di De Zolt}% % lang-it
    Jeśli $Q$ jest figurą zawartą wewnątrz figury $P$ tak, że różnica $P \setminus Q$ ma niepuste wnętrze, to figury $P$ oraz $Q$ nie są równopolowe. % lang-pl
    Se $Q$ è una figura contenuta all'interno di una figura $P$ in modo tale che la differenza $P \setminus Q$ abbia interno non vuoto, allora le figure $P$ e $Q$ non sono equivalenti. % lang-it
\end{axiom}

Starożytni oczywiście nie będą znać teorii miary i postąpią inaczej: wykorzystają zapomniane przez ich potomków rozcięcia. % dissections. % lang-pl
Gli antichi, ovviamente, non conosceranno la teoria della misura e procederanno diversamente: utilizzeranno le dissezioni, dimenticate dai loro posteri. % lang-it
Technika ta pojawi się wielokrotnie, na przykład w dowodzie (I.35), że równoległoboki o~tej samej podstawie i równych wysokościach mają równe pola. % lang-pl
Questa tecnica comparirà più volte, per esempio nella dimostrazione (I.35) che i parallelogrammi con la stessa base e altezze uguali hanno aree uguali. % lang-it
W 1873 roku Henry Perigal (a w 1917 inny Henry, Dudeney) poda dowód twierdzenia Pitagorasa wykorzystujący rozcięcia. % lang-pl
Nel 1873 Henry Perigal (e nel 1917 un altro Henry, Dudeney) fornirà una dimostrazione del teorema di Pitagora basata sulle dissezioni. % lang-it
\index[persons]{Dudeney, Henry}%
\index[persons]{Perigal, Henry}%
Dokładnie opiszą to Eves \cite[s. 194-239]{eves1_1972} i (z nieco inną notacją) Hartshorne \cite[s. 212-221]{hartshorne_2000}. % lang-pl
Lo descriveranno in dettaglio Eves \cite[p. 194-239]{eves1_1972} e (con una notazione leggermente diversa) Hartshorne \cite[p. 212-221]{hartshorne_2000}. % lang-it

Niech $P, Q$ będą jakimiś figurami o tej samej powierzchni. % lang-pl
Siano $P, Q$ due figure di area uguale. % lang-it
Będziemy mówić, że są równopolowe (\emph{equivalent}) i pisać $P \simeq Q$. % lang-pl
Diremo che sono equivalenti (\emph{equivalent}) e scriveremo $P \simeq Q$. % lang-it
Oznaczenie $P \cong Q$ zachowujemy dla zwykłego przystawania. % lang-pl
Riserviamo la notazione $P \cong Q$ per la consueta congruenza. % lang-it

\begin{definition}
    \label{definicja_rownowazne_przez_dodawanie}%
    Dwa wielokąty $P$, $Q$ nazywamy równoważnymi przez rozcinanie, jeśli można rozciąć każdy na skończoną liczbę wielokątnych kawałków o rozłącznych wnętrzach:   % lang-pl
    Due poligoni $P$, $Q$ si dicono equiscomponibili se si può ritagliare ciascuno di essi in un numero finito di pezzi poligonali con interni disgiunti: % lang-it
    \begin{equation}
        P = \bigcup_{k=1}^n P_k, \quad\quad\quad Q = \bigcup_{k=1}^n Q_k,
    \end{equation}
    że kawałek $P_k$ przystaje do $Q_k$ (istnieje pewna izometria $f_k$ taka, że $f_k[P_k] = Q_k$). % lang-pl
    tali che il pezzo $P_k$ sia congruente a $Q_k$ (esiste cioè un'isometria $f_k$ tale che $f_k[P_k] = Q_k$). % lang-it
    Oznaczamy to $P \cong_+ Q$. % lang-pl
    Lo denotiamo con $P \cong_+ Q$. % lang-it
\end{definition}

Eves będzie mówić o przystawaniu przez dodawanie, ale w użyciu będą inne określenia, takie jak \emph{equidecomposable} albo \emph{scissors-congruent}. % lang-pl
Eves parlerà di congruenza per addizione, ma saranno in uso anche altri termini, come \emph{equidecomposable} oppure \emph{scissors-congruent}. % lang-it
Rozcinanie wielokątów będzie przewijać się przez wiele artykułów w Delcie: % lang-pl
Le dissezioni di poligoni ricorreranno in molti articoli su Delta: % lang-it
Joanny Jaszuńskiej w $\Delta_{15}^{8}$, % {https://www.deltami.edu.pl/2015/08/nozyczki-matematyczne/} % lang-pl
di Joanna Jaszuńska in $\Delta_{15}^{8}$, % lang-it
Marka Kordosa w $\Delta_{17}^4$, % {https://www.deltami.edu.pl/2017/04/dedukcja-lokalna-na-przykladzie/} % lang-pl
di Marek Kordos in $\Delta_{17}^4$, % lang-it
Bartłomieja Bzdęgi w $\Delta_{20}^6$ (rozcinanie na równoległoboki), % \todofoot{https://www.deltami.edu.pl/2020/06/rownoleglobok/ zad. 8} % lang-pl
di Bartłomiej Bzdęga in $\Delta_{20}^6$ (dissezione in parallelogrammi), % lang-it
Miłosza Kwiatkowskiego w $\Delta_{25}^3$. % {https://www.deltami.edu.pl/2025/03/wycinanki-i-ukladanki/} % lang-pl
di Miłosz Kwiatkowski in $\Delta_{25}^3$. % lang-it

\begin{definition}
    Dwa wielokąty nazywamy równoważnymi przez doklejanie, jeśli możemy dołączyć do nich pary przystających kawałków tak, by końcowe figury były przystające przez dodawanie. % lang-pl
    Due poligoni si dicono equicompletabili se a ciascuno di essi possiamo aggiungere coppie di pezzi congruenti in modo che le figure finali risultino equiscomponibili. % lang-it
    Formalnie: istnieją przystające przez dodawanie wielokąty $R \cong_+ S$ oraz pewne wielokąty $U_1, U_2, \ldots, U_n$ i $V_1, V_2, \ldots, V_n$ takie, że: $U_1 \cong V_i$, $R = P \sqcup U_1 \sqcup \ldots \sqcup U_n$ oraz $S = Q \sqcup V_1 \sqcup \ldots \sqcup V_n$. % lang-pl
    Formalmente: esistono poligoni equiscomponibili $R \cong_+ S$ e alcuni poligoni $U_1, U_2, \ldots, U_n$ e $V_1, V_2, \ldots, V_n$ tali che: $U_1 \cong V_i$, $R = P \sqcup U_1 \sqcup \ldots \sqcup U_n$ e $S = Q \sqcup V_1 \sqcup \ldots \sqcup V_n$. % lang-it
    Oznaczamy to $P \cong_- Q$. % lang-pl
    Lo denotiamo con $P \cong_- Q$. % lang-it
\end{definition}

\begin{proposition}
    \label{congruent_addition_subtraction}%
    Prawdziwe są implikacje: % lang-pl
    Valgono le seguenti implicazioni: % lang-it
    \begin{itemize}
        \item jeśli $P \cong Q$, to $P \cong_+ Q$. % lang-pl
        \item se $P \cong Q$, allora $P \cong_+ Q$. % lang-it
        \item jeśli $P \cong_+ Q$, to $P \cong_- Q$. % lang-pl
        \item se $P \cong_+ Q$, allora $P \cong_- Q$. % lang-it
        \item jeśli $P \cong_- Q$, to $P$ oraz $Q$ są równopolowe. % lang-pl
        \item se $P \cong_- Q$, allora $P$ e $Q$ sono equivalenti. % lang-it
    \end{itemize}
\end{proposition}

Konieczność dowodu tego twierdzenia dostrzeże W. H. Jackson w 1912 roku, szczegóły można znaleźć u Evesa \cite[s. 198, 199]{eves1_1972} (gdzie pokazuje też, że relacje $\cong_+$ i $\cong_-$ są przechodnie). % lang-pl
La necessità di dimostrare questo fatto sarà notata da W. H. Jackson nel 1912; i dettagli si possono trovare in Eves \cite[p. 198, 199]{eves1_1972} (dove si mostra anche che le relazioni $\cong_+$ e $\cong_-$ sono transitive). % lang-it

Eves wprowadza pojęcie wierzchołka rzutującego (z angielskiego \emph{projecting vertex}, czyli taki, który można oddzielić prostą od pozostałych wierzchołków) i pokazuje, że każdy wielokąt ma co najmniej trzy niewspółliniowe wierzchołki rzutujące. % lang-pl
Eves introduce il concetto di vertice proiettante (dall'inglese \emph{projecting vertex}, cioè un vertice che si può separare dagli altri vertici mediante una retta) e mostra che ogni poligono ha almeno tre vertici proiettanti non allineati. % lang-it
Jest to potrzebne do uzasadnienia, że: % lang-pl
Ciò è necessario per giustificare che: % lang-it

\begin{proposition}
    Każdy wielokąt posiada triangulację: rozcięcie na trójkąty, których wierzchołki są wierzchołkami wyjściowego wielokąta. % lang-pl
    Ogni poligono possiede una triangolazione: una dissezione in triangoli i cui vertici sono vertici del poligono di partenza. % lang-it
\end{proposition}

W przypadku wielokątów wypukłych możemy żądać nawet, by wszystkie trójkąty miały wspólny wierzchołek. % lang-pl
Nel caso di poligoni convessi possiamo persino richiedere che tutti i triangoli abbiano un vertice comune. % lang-it

\begin{proposition}
    Każdy trójkąt jest równoważny przez rozcinanie z równopolowym prostokątem o boku takim samym jak najdłuższy bok trójkąta. % lang-pl
    Ogni triangolo è equiscomponibile con un rettangolo di area uguale avente un lato uguale al lato più lungo del triangolo. % lang-it
\end{proposition}

\begin{corollary}
    Każdy prostokąt jest równoważny przez rozcinanie z kwadratem o takim samym polu. % lang-pl
    Ogni rettangolo è equiscomponibile con un quadrato della stessa area. % lang-it
\end{corollary}

\begin{corollary}
    Każdy wielokąt jest równoważny przez rozcinanie z kwadratem o takim samym polu. % lang-pl
    Ogni poligono è equiscomponibile con un quadrato della stessa area. % lang-it
\end{corollary}

Doszliśmy wreszcie do wyniku Williama Wallace'a (z 1807 roku), znalezionego niezależnie przez Farkasa Bolyaia (1833) i Paula Gerwiena (1835): % lang-pl
Siamo infine giunti al risultato di William Wallace (del 1807), trovato indipendentemente da Farkas Bolyai (1833) e Paul Gerwien (1835): % lang-it
% P. Gerwien, Zerschneidung jeder beliebigen Anzahl von gleichen geradlinigen Figuren in dieselben Stu¨ cke, Journal fu¨ r die reine und angewandte Mathematik 10(1833), 228–234.

% W. Bolyai de Bolya, Tentamen. Iuventutem studiosam in ele- menta matheseos purae, elementaris ac sublı`mioris, methodo intuitiva, evidentiaque huic propria, introducendi. Cum appendice triplici, Maros Va´ sa´ rhelyini; tomus primus 1932, tomus secudus 1933.

\begin{theorem}
[Wallace'a-Bolyaia-Gerwiena] % lang-pl
[di Wallace-Bolyai-Gerwien] % lang-it
\label{twierdzenie_wallace_bolyai_gerwien}%
    Następujące warunki są równoważne: % lang-pl
    Le seguenti condizioni sono equivalenti: % lang-it
    \begin{itemize}
    \item dwie figury są równopolowe, % lang-pl
    \item due figure sono equivalenti, % lang-it
    \item dwie figury są równoważne przez rozcinanie. % lang-pl
    \item due figure sono equiscomponibili. % lang-it
    \end{itemize}
    Punkty rozcięcia można wyznaczyć przy pomocy samego cyrkla i linijki. % lang-pl
    I punti di taglio si possono determinare usando solo compasso e riga. % lang-it
\end{theorem}
% https://en.wikipedia.org/wiki/Wallace-Bolyai-Gerwien_theorem

Laczkovich wzmocni później ten wynik: każdy wielokąt można rozciąć na kawałki i złożyć z nich równopolowy kwadrat tak, by kawałki były przesuwane, ale nie obracane. % lang-pl
Laczkovich rafforzerà in seguito questo risultato: ogni poligono può essere tagliato in pezzi e ricomposto in un quadrato equivalente in modo tale che i pezzi vengano traslati, ma non ruotati. % lang-it
% M. Laczkovich, Equidecomposability and discrepancy: a solution to Tarski’s circle squaring problem, Journal fu¨ r die reine und angewandte Mathematik 404(1990), 77–117.
% M. Laczkovich, Paradoxical decompositions: a survey of recent results, Proc. First European Congress of Mathematics, Vol. II (Paris, 1992), Progress in Mathematics 120, Birkha¨ user, Basel 1994

\subsection{Miara Lebesgue'a} % lang-pl
\subsection{La misura di Lebesgue} % lang-it

Skończymy na odległych czasach. % lang-pl
Concluderemo con tempi lontani. % lang-it
Tysiące lat po Euklidesie pojęcie rozmiaru (nie tylko długości, ale przede wszystkim pola figur płaskich oraz objętości brył) sformalizuje się przez miarę Jordana, albo jej uogólnienie --- miarę Lebesgue'a. % lang-pl
Migliaia di anni dopo Euclide, il concetto di grandezza (non solo la lunghezza, ma soprattutto l'area delle figure piane e il volume dei solidi) si formalizzerà attraverso la misura di Jordan, oppure la sua generalizzazione: la misura di Lebesgue. % lang-it
O mierze wspomina Hartshorne \cite[s. 205]{hartshorne_2000}; jest to funkcja, która przyporządkowuje każdej figurze $P$ pewną nieujemną liczbę $[P]$ (zwaną polem) taką, że: % lang-pl
Della misura parla Hartshorne \cite[p. 205]{hartshorne_2000}; si tratta di una funzione che associa a ogni figura $P$ un numero non negativo $[P]$ (detto area) tale che: % lang-it
\begin{itemize}
    \item pole każdego trójkąta jest dodatnie: $[\triangle] > 0$; % lang-pl
    \item l'area di ogni triangolo è positiva: $[\triangle] > 0$; % lang-it
    \item przystające trójkąty mają równe pola; % lang-pl
    \item triangoli congruenti hanno aree uguali; % lang-it
    \item jeśli dwie figury $P, Q$ nie nachodzą na siebie, to $[P \cup Q] = [P] + [Q]$. % lang-pl
    \item se due figure $P, Q$ non si sovrappongono, allora $[P \cup Q] = [P] + [Q]$. % lang-it
\end{itemize}

Z istnienia takiej funkcji wynika aksjomat (Z), wysłowiony przez nas niżej. % lang-pl
Dall'esistenza di una tale funzione segue l'assioma (Z), da noi enunciato più sotto. % lang-it
Guzicki \cite[s. 38-68]{guzicki_2021} również wspomni o mierze Jordana. % lang-pl
Anche Guzicki \cite[p. 38-68]{guzicki_2021} menzionerà la misura di Jordan. % lang-it
Zauważy jednak, że w polskich szkołach każdą figurę można rozciąć na wielokąty i wycinki kołowe, co prowadzi do prostszej aksjoamtyzacji (i naturalnie do dowodu twierdzenia \ref{twierdzenie_wallace_bolyai_gerwien}, Wallace'a-Bolyaia-Gerwiena). % lang-pl
Osserverà però che, nelle scuole polacche, ogni figura può essere ritagliata in poligoni e settori circolari, il che porta a un'assiomatizzazione più semplice (e naturalmente alla dimostrazione del teorema \ref{twierdzenie_wallace_bolyai_gerwien}, di Wallace-Bolyai-Gerwien). % lang-it

Pole jest wyznaczone jednoznacznie z dokładnością do stałej. % lang-pl
L'area è determinata in modo univoco a meno di una costante. % lang-it
Wartość stałej ustala się zazwyczaj przez przyjęcie, że trójkąt o podstawie $a$ oraz wysokości $h$ ma pole $\frac 1 2 ah$. % lang-pl
Il valore della costante si fissa di solito assumendo che un triangolo di base $a$ e altezza $h$ abbia area $\frac 1 2 ah$. % lang-it

\subsection{Zasada Cavalieriego} % lang-pl
\subsection{Il principio di Cavalieri} % lang-it

Sekcję zamkniemy ciekawostką związaną z zasadą Cavalieriego. % lang-pl
Chiuderemo la sezione con una curiosità legata al principio di Cavalieri. % lang-it
Eves pokaże w 1991 roku, że każde dwa równopolowe trójkąty można umieścić na płaszczyźnie tak, by każda prosta równoległa do danej przecinała obydwa na odcinku o tej samej długości, wspomni o tym Jarosław Górnicki w $\Delta_{11}^{11}$. % lang-pl
Eves mostrerà nel 1991 che due triangoli equivalenti qualsiasi si possono sempre collocare nel piano in modo che ogni retta parallela a una retta data intersechi entrambi in un segmento della stessa lunghezza; ne parlerà Jarosław Górnicki in $\Delta_{11}^{11}$. % lang-it
% {https://www.deltami.edu.pl/2011/11/osobliwosc-trojkatow/}

%

%

\pol{\section{Origami}}
\ita{\section{Origami}}
\label{section_origami}
\index{origami}

\pol{Aksjomaty Huzity-Justina (albo Huzity-Hatoriego) to zestaw siedmiu aksjomatów, które opisują wszystkie możliwe pojedyncze zgięcia w geometrii origami na płaszczyźnie.}
\ita{Gli assiomi di Huzita-Justin (o di Huzita-Hatori) sono un insieme di sette assiomi che descrivono tutte le possibili pieghe singole nella geometria origami sul piano.}
\pol{\index{aksjomat!Huzity-Hatoriego}}\ita{\index{assioma!di Huzita-Hatori}}%
\pol{\index{aksjomat!Huzity-Justina}}\ita{\index{assioma!di Huzita-Justin}}%
\pol{Niechf $p_1, p_2$ będą dwoma (różnymi) punktami, zaś $l_1, l_2$ dwoma różnymi prostymi.}
\ita{Siano $p_1, p_2$ due punti (distinti), e $l_1, l_2$ due rette distinte.}

\begin{enumerate}
    \item \pol{Istnieje dokładnie jedno zgięcie, które przechodzi przez punkty $p_1$ i $p_2$.}
    \ita{Esiste esattamente una piega che passa per i punti $p_1$ e $p_2$.}
    \item \pol{Istnieje dokładnie jedno zgięcie, które przenosi punkt $p_1$ na $p_2$.}
    \ita{Esiste esattamente una piega che porta il punto $p_1$ su $p_2$.}
    \item \pol{Istnieje zgięcie, które przenosi prostą $l_1$ na $l_2$.}
    \ita{Esiste una piega che porta la retta $l_1$ su $l_2$.}
    \item \pol{Istnieje dokładnie jedno zgięcie przez $p_1$, które jest prostopadłe do $l_1$.}
    \ita{Esiste esattamente una piega passante per $p_1$ che è perpendicolare a $l_1$.}
    \item \pol{Istnieje zgięcie przez $p_2$, które przenosi punkt $p_1$ na prostą $p_2$.}
    \ita{Esiste una piega passante per $p_2$ che porta il punto $p_1$ sulla retta $p_2$.}
    \item \pol{Istnieje zgięcie, które przenosi $p_1$ na $l_1$, zaś $p_2$ na $l_2$.}
    \ita{Esiste una piega che porta $p_1$ su $l_1$ e $p_2$ su $l_2$.}
    \item \pol{Istnieje zgięcie prostopadłe do $l_2$, które przenosi $p_1$ na $l_2$.}
    \ita{Esiste una piega perpendicolare a $l_2$ che porta $p_1$ su $l_2$.}
\end{enumerate}

\pol{Zestaw nie jest minimalny, ale wyczerpuje wszystkie możliwe ruchy.}
\ita{L'insieme non è minimale, ma esaurisce tutti i movimenti possibili.}
\pol{Wszystkie znajdzie francuski składacz Jacques Justin \cite{justin_1986} w 1986 roku.}
\ita{Tutti saranno trovati dal piegatore francese Jacques Justin \cite{justin_1986} nel 1986.}
\index[persons]{Justin, Jacques}%
\pol{Później część z nich będzie odkrywana na nowo: przez japońsko-włoskiego matematyka Humiakiego Huzitę (od 1 do 6, około 1991 roku), przez Davida Aucklyego i Johna Clevelanda \cite{cleveland_1995} (od 1 do 5 w 1995 roku), Koshiro Hatoriego (aksjomat nr 7 w 2001 roku) i Roberta Langa \cite{alperin_2009} (także aksjomat nr 7).}
\ita{In seguito alcuni di essi saranno riscoperti: dal matematico nippo-italiano Humiaki Huzita (dall'1 al 6, intorno al 1991), da David Auckly e John Cleveland \cite{cleveland_1995} (dall'1 al 5 nel 1995), da Koshiro Hatori (assioma n. 7 nel 2001) e da Robert Lang \cite{alperin_2009} (anch'egli l'assioma n. 7).}
\index[persons]{Huzita, Humiaki}%
\index[persons]{Auckly, David}%
\index[persons]{Cleveland, John}%
\index[persons]{Hatori, Koshiro}%
\index[persons]{Lang, Robert}%
% TODO: Alperin =  Roger C. Alperin; Robert J. Lang (2009). "One-, Two-, and Multi-Fold Origami Axioms" (PDF). 4OSME. A K Peters. Archived from the original (PDF) on 2022-02-13. Retrieved 2012-04-20.

\pol{Aksjomat nr 2 znajduje symetralną odcinka $p_1 p_2$; nr 3 znajduje dwusieczną kąta, na ramionach którego leżą proste $l_1$, $l_2$.}
\ita{L'assioma n. 2 trova l'asse del segmento $p_1 p_2$; il n. 3 trova la bisettrice dell'angolo sui cui lati giacciono le rette $l_1$, $l_2$.}
\pol{Może to być nieoczywiste, ale aksjomat nr 5 jest równoważny przecięciu okręgu prostej.}
\ita{Può non essere evidente, ma l'assioma n. 5 è equivalente all'intersezione di un cerchio e una retta.}
\pol{Aksjomat nr 6 znajduje prostą styczną do dwóch parabol.}
\ita{L'assioma n. 6 trova una retta tangente a due parabole.}
\pol{(To zgięcie nazywa się zgięciem Beloch, ponieważ Margharita Beloch pokaże w 1936, że jest to równoważne co rozwiązaniu wielomianowych równań trzeciego stopnia).}
\ita{(Questa piega si chiama piega di Beloch, perché Margharita Beloch mostrerà nel 1936 che essa è equivalente alla risoluzione di equazioni polinomiali di terzo grado).}
\index[persons]{Beloch, Margharita}%
% https://en.wikipedia.org/wiki/Mathematics_of_paper_folding

\pol{Aksjomaty od pierwszego do czwartego są słabsze niż cyrkiel i linijka: pozwalają zbudować liczby pitagorejskie, najmniejsze ciało zawierające wszystkie liczby wymierne, które oprócz (dowolnej) liczby $x$ zawierają też $\sqrt{x^2 + 1}$.}
\ita{Gli assiomi dal primo al quarto sono più deboli di riga e compasso: permettono di costruire i numeri pitagorici, il più piccolo campo contenente tutti i numeri razionali che, oltre a un numero (qualunque) $x$, contiene anche $\sqrt{x^2 + 1}$.}
\pol{Aksjomaty od 1 do 5 są dokładnie tak samo mocne jak cyrkiel i linijka, nr 6 pozwala na rozwiązanie problemów delfickiej wyroczni, zaś nr 7 nie daje nic nowego.}
\ita{Gli assiomi da 1 a 5 sono esattamente potenti quanto riga e compasso, il n. 6 permette di risolvere i problemi dell'oracolo di Delfi, mentre il n. 7 non dà nulla di nuovo.}
\pol{Ale zgięcie, które opisuje, jest możliwe do wykonania na kartce papieru, więc aksjomat nie jest całkiem niepotrzebny.}
\ita{Ma la piega che descrive è possibile da eseguire su un foglio di carta, quindi l'assioma non è del tutto inutile.}

\pol{Jorge Lucero w 2017 dołoży ósmy aksjomat, że istnieje zgięcie wzdłuż prostej $l_1$.}
\ita{Jorge Lucero nel 2017 aggiungerà un ottavo assioma, secondo cui esiste una piega lungo la retta $l_1$.}
\index[persons]{Lucero, Jorge}%

\pol{Osiem zadań z konkursu matematycznego origami ,,Żuraw'' znajdziemy w $\Delta_{12}^5$.}
\ita{Otto problemi del concorso matematico di origami ,,Żuraw'' si trovano in $\Delta_{12}^5$.}
\pol{Dwie studentki, Barbara Ciesielska i Agnieszka Kowalczyk, opiszą w $\Delta_{16}^7$ konstrukcję odcinka długości $\sqrt[3]{2}$ oraz trysekcję kąta przy użyciu origami.}
\ita{Due studentesse, Barbara Ciesielska e Agnieszka Kowalczyk, descriveranno in $\Delta_{16}^7$ la costruzione di un segmento di lunghezza $\sqrt[3]{2}$ e la trisezione dell'angolo mediante origami.}

\pol{Todo: Hartshorne \cite[s. 249]{hartshorne2000}.}
\ita{Todo: Hartshorne \cite[p. 249]{hartshorne2000}.}

%

%